% Papel milimetrado (A4) — Eric Lucero González — https://ericlucero.dev
% Licencia CC BY-NC-SA 4.0. Compilar con pdfLaTeX (dos pasadas: latexmk lo hace solo).
\documentclass[a4paper]{article}
\usepackage[T1]{fontenc}
\usepackage[spanish]{babel}
\usepackage{lmodern}
\usepackage{tikz}
\pagestyle{empty}

% ── Configuración (mm) ─────────────────────────────────────
\def\anchoPapel{210}
\def\altoPapel{297}
\def\margenMin{10}   % margen mínimo alrededor de la rejilla

\begin{document}
\begin{tikzpicture}[
    remember picture, overlay, shift={(current page.south west)},
    x=1mm, y=1mm,
    mm lines/.style={orange!35, ultra thin},
    half cm lines/.style={orange!55, very thin},
    cm lines/.style={orange!80!black, thin}
  ]
  % Centímetros enteros que caben; la rejilla queda centrada en la hoja
  \pgfmathtruncatemacro{\nx}{(\anchoPapel - 2*\margenMin) / 10}
  \pgfmathtruncatemacro{\ny}{(\altoPapel  - 2*\margenMin) / 10}
  \pgfmathsetmacro{\ancho}{\nx * 10}
  \pgfmathsetmacro{\alto}{\ny * 10}
  \pgfmathsetmacro{\x}{(\anchoPapel - \ancho) / 2}
  \pgfmathsetmacro{\y}{(\altoPapel  - \alto) / 2}

  \begin{scope}[shift={(\x,\y)}]
    \draw[step=1,  mm lines]      (0,0) grid (\ancho,\alto);
    \draw[step=5,  half cm lines] (0,0) grid (\ancho,\alto);
    \draw[step=10, cm lines]      (0,0) grid (\ancho,\alto);
  \end{scope}
\end{tikzpicture}
\end{document}
